% AIDED: Assisted Analysis of the Dynamic Electron Density
% Copyright (C) 2025 J. Robert Michael, PhD. All Rights Reserved.
\documentclass{article}

% ---------- generic packages ----------
\usepackage[utf8]{inputenc}   % << handles any UTF-8 that sneaks in
\usepackage[T1]{fontenc}
\usepackage{amsmath,amssymb,bm}
\usepackage{hyperref}
\usepackage{listings}         % nicer verbatim for the Python code
\usepackage{aided_preamble}   % your own colours/macros, if any
\usepackage[numbers,sort&compress]{natbib}

% ---------- canonical notation ----------
% (defined once – used everywhere)
\newcommand*\expa{\alpha}
\newcommand*\expb{\beta}
\newcommand*\expab{\gamma}           % γ = α + β
\newcommand*\prefac{N_{\expa\expb}}  % √[γ⁻³ det W]

\newcommand*\veca{\bm a}
\newcommand*\vecb{\bm b}
\newcommand*\vecc{\bm c}

\newcommand*\matG{\bm G}             % γ I₃
\newcommand*\matU{\bm U}             % vibration tensor
\newcommand*\matW{\bm W}             % (G⁻¹ + U)⁻¹

\newcommand*\RR{\mathbb R}
\newcommand*\dx[1]{\partial_{#1}}

% --------------------------------------------------------------------------
% listings setup
\lstset{
  basicstyle=\ttfamily\small,
  numbers=left, numberstyle=\tiny, numbersep=4pt,
  frame=single, framerule=0.4pt,
  breaklines=true, breakindent=1em
}

\title{\textbf{AIDED}\\\large Assisted Analysis of the Dynamic Electron Density}
\author{J. Robert Michael, PhD}
\date{\today}

\begin{document}\maketitle
\tableofcontents

%---------------------------------------------------------------------------
\section{Introduction}

\emph{AIDED} evaluates both the static and the vibration-averaged
(“dynamic”) electron density of molecular systems, together with derived
topological properties such as critical points and zero-flux surfaces.  A
key design goal is to \emph{reuse the primitive-Gaussian basis} exported by
standard quantum-chemistry programs so that no extra integrals are needed.

\bigskip

%---------------------------------------------------------------------------
\section{Static electron density}
\label{sec:static}

The (static) electron density $\rho(\bm r)$ of a system with $N$
electrons is\footnote{Spin is omitted for clarity.}
\[
  \rho(\bm r)=
  N\int
    \Psi(\bm r,\bm r_2,\dots,\bm r_N;\bm R)\,
    \Psi^{*}(\bm r,\bm r_2,\dots,\bm r_N;\bm R)\,
    d\bm r_2\cdots d\bm r_N ,
\]
where $\Psi$ is the ground-state wave-function at the fixed nuclear
geometry $\bm R$.

%...........................................................................
\subsection{Basis expansion}

Each \emph{atomic} orbital (AO) is a contraction of primitive Gaussians
$g_k$:
\begin{equation}\label{eq:phi_j}
  \phi_j(\bm r)=\sum_{k} g_{jk}\,g_k(\bm r).
\end{equation}
A molecular orbital (MO) is an AO-linear combination,
\begin{equation}\label{eq:Psi_i}
  \Psi_i(\bm r)=\sum_{\mu} c_{i\mu}\,\phi_\mu(\bm r)
              =\sum_{k} C_{ik}\,g_k(\bm r),
  \qquad C_{ik}=\sum_{\mu}c_{i\mu}\,g_{\mu k}.
\end{equation}

The AO-resolved one-particle density matrix is
\[
  P_{\mu\nu}=2\sum_i n_i\,c_{i\mu}c_{i\nu},
\]
with $n_i\in\{0,1,2\}$ the MO occupations.  
Inserting \eqref{eq:phi_j} gives the primitive-basis density matrix
\begin{equation}\label{eq:Dkl}
  D_{kl}=\sum_{\mu,\nu} g_{\mu k}\,P_{\mu\nu}\,g_{\nu l}
        =\sum_{i} n_i\,C_{ik}C_{il},
\end{equation}
so that the density at any grid point is the simple quadratic form
\[
  \rho(\bm r)=\sum_{k,l} D_{kl}\,g_k(\bm r)\,g_l(\bm r)
            =\bm g^{\!\mathsf T}(\bm r)\,D\,\bm g(\bm r).
\]

The \textsc{AIDED} code builds~$D$ once from the $\mathtt{*.wfn}$ MO
coefficients, then evaluates $\rho(\bm r)$ on the desired grid.

%---------------------------------------------------------------------------
\section{Mean-square displacement amplitudes (MSDA)}
\label{sec:msda}

Normal-mode analysis of a molecule with $M$ nuclei gives a $(3M)$-variate
Gaussian probability distribution for nuclear displacements
$\bm u=\bm R-\bm R^{0}$:
\begin{equation}\label{eq:Pu}
  P(\bm u)=\frac{\exp\!\bigl(-\tfrac12\bm u^{\mathsf T}\matU^{-1}\bm u\bigr)}
                {(2\pi)^{3M/2}\sqrt{\det\matU}},
\end{equation}
where $\matU=\langle\bm u\bm u^{\mathsf T}\rangle_T$ is the Cartesian MSDA
matrix.  Diagonalising the mass-weighted Hessian yields
\[
  \matU=\bm L\,\bm\Delta\,\bm L^{\mathsf T},
  \quad
  \Delta_{jj}=\frac{h}{8\pi^{2}\nu_j}
              \coth\!\bigl(\tfrac{h\nu_j}{2k_{\mathrm B}T}\bigr),
\]
implemented verbatim in
\href{https://github.com/drjrm3/aided/blob/main/aided/io/vib/gaussian.py}
{\texttt{gaussian.py::gen\_msda}}.

%---------------------------------------------------------------------------
\section{Dynamic electron density}\label{sec:dynamic}

\subsection{Primitive $s$–$s$ unit}

\[
  \expab=\expa+\expb,\qquad
  \vecc=\frac{\expa\veca+\expb\vecb}{\expab},\qquad
  \matG=\expab I_{3},\qquad
  \matW=(\matG^{-1}+\matU)^{-1}.
\]

\begin{equation}
  \boxed{%
    g_{ss}^{\text{dyn}}(\bm x)=
      \prefac
      \exp\!\Bigl[
        -\tfrac12(\bm x-\vecc)^{\!\mathsf T}\matW^{-1}(\bm x-\vecc)
        -\tfrac12\frac{\expa\expb}{\expab}\lvert\veca-\vecb\rvert^{2}
      \Bigr]},
  \quad
  \prefac=\bigl[\expab^{-3}\det\matW\bigr]^{1/2}.
  \label{eq:g_dyn_ss}
\end{equation}

%----------------------------------------------------------------
\subsection{Primitive $s$--$p$ unit}\label{sec:sp_unit}

The Cartesian $p$ factor is generated by differentiating
$g_{ss}^{\text{dyn}}$ with respect to one nuclear centre
(choose $\vecb$ without loss of generality).  Because
\[
  \dx{b_\nu}\,\vecc =\frac{\expb}{\expab}\,\bm e_\nu,
  \qquad
  \dx{b_\nu}\,|\veca-\vecb|^{2} =-2(\veca_\nu-\vecb_\nu),
\]
one obtains
\begin{align}
  g_{sp_\nu}^{\text{dyn}}(\bm x)
  &\equiv -\frac{1}{2\expb}\,\dx{b_\nu}g_{ss}^{\text{dyn}}(\bm x) \\[4pt]
  &=\underbrace{\bigl[
      \tfrac{\expb}{\expab}( \matW^{-1}(\bm x-\vecc) )_\nu
      \;+\;\tfrac{\expa\expb}{\expab}\,(\veca_\nu-\vecb_\nu)
    \bigr]}_{\displaystyle L_\nu(\bm x)}
    \;g_{ss}^{\text{dyn}}(\bm x).
  \label{eq:g_dyn_sp}
\end{align}

\paragraph{Linear prefactor.}
The bracketed quantity
\[
  L_\nu(\bm x)=
    \frac{\expb}{\expab}\bigl(\matW^{-1}(\bm x-\vecc)\bigr)_\nu
    +\frac{\expa\expb}{\expab}\,(\veca_\nu-\vecb_\nu)
\]
is exactly the “$C_{4},C_{5},C_{6}$” coefficient that the
\textsc{Fortran} and NumPy kernels store for $p_x,p_y,p_z$ on centre $j$.

(Interchanging $\expa\leftrightarrow\expb$ and
$\veca\leftrightarrow\vecb$ gives the $p$ factor on centre $i$).

%----------------------------------------------------------------
\subsection{General Cartesian exponents}\label{sec:general_cartesian}

For arbitrary Cartesian angular momenta
$\bm L_i\!=\!(\ell_i,m_i,n_i)$ and $\bm L_j\!=\!(\ell_j,m_j,n_j)$,
repeat the differentiation:

\[
  g_{\bm L_i;\bm L_j}^{\text{dyn}}(\bm x)=
    \prod_{\mu=1}^{3}\bigl(\dx{a_\mu}\bigr)^{\ell_i^{(\mu)}}
    \prod_{\nu=1}^{3}\bigl(\dx{b_\nu}\bigr)^{\ell_j^{(\nu)}}
    \,g_{ss}^{\text{dyn}}(\bm x).
\]

All such derivatives factorise into
\[
  g_{\bm L_i;\bm L_j}^{\text{dyn}}(\bm x)=
     g_{ss}^{\text{dyn}}(\bm x)\,
     \prod_{\mu=1}^{3}\bigl[C_{\mu}(\bm x)\bigr]^{\ell_i^{(\mu)}}
     \prod_{\nu=1}^{3}\bigl[C_{3+\nu}(\bm x)\bigr]^{\ell_j^{(\nu)}},
\]
where the six \emph{linear} terms are
\[
  C_{1,2,3}(\bm x)=
    \frac{\expa}{\expab}\,\matW^{-1}(\bm x-\vecc)
    -\frac{\expa\expb}{\expab}\,(\veca-\vecb),
\qquad
  C_{4,5,6}(\bm x)=
    \frac{\expb}{\expab}\,\matW^{-1}(\bm x-\vecc)
    +\frac{\expa\expb}{\expab}\,(\veca-\vecb),
\]
\emph{i.e.}\ the $p$-type prefactors on centres $i$ and $j$.
Higher-$l$ Gaussians follow by simply raising these linear factors to the
required powers—the same rule implemented in
`gen\_gcon` and in the NumPy kernel of
Listing \ref{lst:python}.


%----------------------------------------------------------------
\subsection{Reference NumPy implementation}

\begin{lstlisting}[caption={Vectorised NumPy kernel matching Eq. (10)},
                   label={lst:python}]
import numpy as np

def dynamic_gaussian(x, A, B, alpha, beta, U,
                     lmn_i=(0,0,0), lmn_j=(0,0,0)):
    gamma = alpha + beta
    C     = (alpha*A + beta*B) / gamma
    W     = np.linalg.inv(np.eye(3)/gamma + U)

    prefactor = (gamma**-3 * np.linalg.det(W))**0.5
    Eg  = np.exp(-0.5*alpha*beta/gamma * np.dot(A-B, A-B))
    expo = np.exp(-0.5 * (x-C) @ np.linalg.inv(W) @ (x-C))
    g_ss = prefactor * Eg * expo

    # Ci factors omitted for brevity ...
    return g_ss
\end{lstlisting}



%================================================================
\section*{NOTES}
This document is based off of my PhD thesis \cite{MichaelPhD2013} and is
based off of Scheringer \cite{Scheringer1976} which is referenced in
\cite{Michael2015JMC}.
and is originally implementedin the DENPROP package \cite{DENPROP}.

This document was generated with the help of AI LLM models with guidance
from my previous work.

%================================================================
\begin{thebibliography}{99}
\bibitem{Scheringer1976}
J.\ Scheringer and H.\ Reitz,
\emph{Direct Calculation of Dynamic Densities},
\textit{Acta Crystallographica A}, \textbf{32} (1976) 271–276.

\bibitem{Michael2015JMC}
J.~R. Michael and T.~Koritsanszky,
\newblock Validation of convolution approximation to the thermal-average
  electron density,
\newblock {\em Journal of Mathematical Chemistry} \textbf{53} (2015) 250--259.

\bibitem{MichaelPhD2013}
J.~R. Michael,
\newblock \emph{Direct Calculation of Dynamic Electron Density:  
         Theoretical Development and Experimental Validation},
\newblock Ph.D.\ thesis, Middle Tennessee State University (2013).

\bibitem{DENPROP}
A.~Volkov,T.~Koritsanszky,M.~Chodkiewicz, and H.~F.~King
\newblock On the basis-set dependence of local and integrated electron-density
  properties: Application of a new computer program for quantum-chemical
  density analysis,
\newblock \emph{Journal of Computational Chemistry} \textbf{25} (2004) 166–183.
  \href{https://doi.org/10.1002/jcc.10434}{doi:10.1002/jcc.10434}.


\end{thebibliography}

\end{document}
